\scanpage{71}
% archive foliation: 70
C'est trivial.

\deleted{Cas particuliers\quad \uline{Proposition} Dans les cas des $\unsure{\mathrm I'}$ ou $\unsure{\mathrm{II}''}$, si $W^a(\mathfrak X)$ existe et}

\uline{Lemme 2} Soient $\mathfrak X_i$ des ouverts de $\mathfrak X/W(S)$, on suppose \deleted{que $\underline W^a(\mathfrak X)$ existe} que \deleted{les $\mathfrak X_i$ dans $\mathfrak X$ \unsure{vérifient} les conditions du Lemme 1 (en particulier} les $W^a(\mathfrak X_i)$ existent \deleted{et que} \deleted{)} et que les $\mathfrak X_i\cap\mathfrak X_j$ dans $\mathfrak X_i$ satisfont les $=$ conditions \inserted{du lemme 1}, [donc $\underline W^a(\mathfrak X_i\cap\mathfrak X_j)$ existe\unsure{\}} et est un ouvert dans $\underline W^a(\mathfrak X_i)$], \inserted{et \unsure{bis-bis|bis-idem} pour les $\mathfrak X_i\cap\mathfrak X_j\cap\mathfrak X_{\unsure{h|k}}$ dans $\mathfrak X_i\cap\mathfrak X_j$}. On suppose

\textcircled{$\beta$} Pour $\forall$ $T/S$, et tout morphisme $\underline W(T)\xrightarrow{f}\mathfrak X$, on a $T=\bigcup_i T_i$ (les $T_i$ ouverts de $T$) tel que
\[
f(W(T_i))\subset\mathfrak X_i
\]
Dans ces conditions, $\underline W^a(\mathfrak X)$ existe, et est la réunion des $\underline W^a(\mathfrak X_i)$, qui s'identifient à des ouverts de $\underline W^a(\mathfrak X)$
